\documentclass[../../main.tex]{subfiles}

\graphicspath{{../../image/}{../../figure/}}

\begin{document}

\section{双曲性}\label{section:ch03-双曲性}

回到第~I 章的设定，我们把 $\mathbb{R}^{k}$ 视作 $\mathbb{R}^{m}\times\mathbb{R}$：其中 $\mathbb{R}^{m}$（$m=k-1$）是"空间"，典型点记为 $x$；$\mathbb{R}$ 是"时间"，典型值记为 $t$，于是 $X=(x,t)$。我们记 $\partial_{t}=\partial/\partial X_{k}$，$\partial_{\alpha}=\partial/\partial X_{\alpha}$，$\alpha=1,\cdots,m$。沿用前面各章的记号，用 $\operatorname{div}$ 表示对 $\mathbb{R}^{m}$ 中 $x$ 变量求散度；在矩阵运算中 $\operatorname{div}$ 作用于行向量。我们还沿用记号~1.4.2%

，它在本书中始终有效：$D$ 表示微分 $[\partial/\partial U^{1},\ldots,\partial/\partial U^{n}]$，视为行算子。

把 $G_{k}$ 改记为 $H$，并把符号 $G$ 重新指派为以 $(G_{1},\ldots,G_{m})$ 为列向量的 $n\times m$ 矩阵，于是平衡律方程组\eqref{align:ch01-closed-system}%

可改写为
\begin{align}
\partial_{t}H(U(x,t),x,t)+\operatorname{div}G(U(x,t),x,t)=\Pi(U(x,t),x,t).
\label{align:ch03-平衡律一般形式}
\end{align}

\begin{definition}[沿 $t$ 方向的双曲性]\label{definition:ch03-沿t方向双曲性}
对任意固定的 $U\in\mathcal{O}$、$(x,t)\in\mathscr{X}$ 以及 $\nu\in\mathbb{S}^{m-1}$，若 $n\times n$ 矩阵 $DH(U,x,t)$ 非奇异，且特征值问题
\begin{align}
\left[\sum_{\alpha=1}^{m}\nu_{\alpha}DG_{\alpha}(U,x,t)-\lambda DH(U,x,t)\right]R=0
\label{align:ch03-特征值问题}
\end{align}
有实特征值 $\lambda_{1}(\nu;U,x,t),\cdots,\lambda_{n}(\nu;U,x,t)$（称为\emph{特征速度}）以及 $n$ 个线性无关的特征向量 $R_{1}(\nu;U,x,t),\cdots,R_{n}(\nu;U,x,t)$，则称平衡律方程组\eqref{align:ch03-平衡律一般形式}是\emph{沿 $t$ 方向双曲的}。
\end{definition}

一类极为重要的方程组是\emph{对称双曲}平衡律方程组\eqref{align:ch03-平衡律一般形式}：对任意 $U\in\mathcal{O}$ 与 $(x,t)\in\mathscr{X}$，$n\times n$ 矩阵 $DG_{\alpha}(U,x,t)$（$\alpha=1,\cdots,m$）对称，且 $DH(U,x,t)$ 对称正定。

双曲性的定义可以借助\cref{section:ch01-弱阵面与激波阵面}%

中引入的波前概念作自然的几何解释。平衡律方程组\eqref{align:ch03-平衡律一般形式}的一个波前 $\mathscr{F}$ 可视为一族以 $t$ 为参数的 $m-1$ 维流形，即在空间中传播的一张曲面。在这一图景下，把 $\mathscr{F}$ 的法向量 $N$ 重整化为 $N=(\nu,-s)$，其中 $\nu\in\mathbb{S}^{m-1}$，则波沿方向 $\nu$ 以速度 $s$ 传播。因此，把\eqref{align:ch03-特征值问题}与\eqref{align:ch01-weak-front-jump}%

对照可知：一个含 $n$ 个平衡律的方程组是双曲的，当且仅当在任意空间方向上都允许传播 $n$ 个不同的弱波。\eqref{align:ch03-特征值问题}的特征值确定这些波的传播速度，相应的特征向量则规定其振幅方向。

当 $\mathscr{F}$ 是激波阵面时，\eqref{align:ch01-Rankine-Hugoniot}%

在当前记号下可写为
\begin{align}
-s\big[H(U_{+},x,t)-H(U_{-},x,t)\big]+\big[G(U_{+},x,t)-G(U_{-},x,t)\big]\nu=0,
\label{align:ch03-RH一般形式}
\end{align}
称为\emph{Rankine--Hugoniot 跳跃条件}。由\cref{theorem:ch01-BV解刻画}%

，对于任何在测度意义下满足\eqref{align:ch03-平衡律一般形式}、且属于 $BV_{\mathrm{loc}}$ 类的函数 $U$，这一条件应在其近似跳跃间断的每一点处成立。

显然，双曲性在状态向量的任意变换 $U^{*}=U^{*}(U,x,t)$（其中 $U^{*}(\cdot,x,t)$ 对每一固定 $(x,t)\in\mathscr{X}$ 都是微分同胚）下保持不变。特别地，由于 $DH(U,x,t)$ 非奇异，我们至少可以局部地取 $H$ 作为新的状态向量。于是在本质上不失一般性，可以把研究范围限制为具有如下特殊形式的双曲平衡律方程组：
\begin{align}
\partial_{t}U(x,t)+\operatorname{div}G(U(x,t),x,t)=\Pi(U(x,t),x,t).
\label{align:ch03-正则形式}
\end{align}

为简单方便起见，今后我们把特殊形式\eqref{align:ch03-正则形式}视为\emph{正则形式}。但读者应注意：在处理连续介质物理中导出的平衡律方程组时，沿用自然给出的状态向量可能更为有利，即使这样做必须面对较\eqref{align:ch03-正则形式}更复杂的一般形式\eqref{align:ch03-平衡律一般形式}。

\end{document}
